% sections/02_related.tex
\section{Related Work}
\label{sec:related}

\subsection{LLM-based Unit Test Generation}
\label{subsec:related-llm}

Nhiều nghiên cứu gần đây đã đánh giá khả năng sinh unit test tự động của các mô hình ngôn ngữ lớn (LLM) trên các mã nguồn thực tế. Huang et al.~\cite{huang2026ult} giới thiệu benchmark ULT gồm 3.909 hàm Python (CC $\geq$ 10) được xây dựng để chống rò rỉ dữ liệu (data contamination), và chỉ ra rằng độ chính xác Pass@5 trên benchmark sạch chỉ đạt 12,57\%, thấp hơn nhiều so với 44,45\% trên benchmark cũ bị rò rỉ -- cho thấy các kết quả công bố trước đây có thể bị thổi phồng đáng kể. Wang et al.~\cite{wang2024hits} (HITS) kết hợp GPT-3.5-turbo với kỹ thuật phân rã phương thức (method slicing) và chain-of-thought, đạt line coverage 55,09\% trên 114 phương thức Java phức tạp (CC $>$ 10), vượt trội các baseline LLM khác nhưng vẫn thấp hơn nhiều so với các phương thức có độ phức tạp trung bình được báo cáo trong các nghiên cứu khác. Yuan et al.~\cite{yuan2024nomoretests} đánh giá ChatGPT trên 13 dự án Java mã nguồn mở, đạt tỷ lệ biên dịch 91,2\% và độ bao phủ tương đương test viết thủ công, song không báo cáo mutation score. Schäfer et al.~\cite{schafer2024empirical} so sánh nhiều LLM (GPT-4, Claude 3.5, Command-R, Llama 3.1) trên cả Java và Python, ghi nhận hiệu năng trên Python thấp hơn đáng kể so với Java ở mọi mô hình được thử nghiệm. Riêng với GPT-4o-mini -- mô hình được sử dụng trong nghiên cứu này -- công trình gần đây nhất~\cite{fse2025gpt4ominigranularity} áp dụng kỹ thuật hybrid-granularity prompting trên Defects4J (Java), chỉ đạt line/branch/mutation coverage lần lượt 52,30\%/38,84\%/36,76\%, với khoảng 43\% test sinh ra không pass được.

Điểm chung của nhóm nghiên cứu này là: (1) hiệu năng LLM suy giảm mạnh khi độ phức tạp phương thức tăng (HITS); (2) chưa đến một nửa số công trình báo cáo mutation score bên cạnh coverage, dù coverage cao không đồng nghĩa với khả năng phát hiện lỗi tốt~\cite{huang2026ult}; (3) phần lớn tập trung vào Java, trong khi Python -- ngôn ngữ được dùng trong nghiên cứu này -- cho kết quả thấp và ít được kiểm soát về độ phức tạp hàm một cách hệ thống; và (4) rủi ro data contamination có thể làm sai lệch đáng kể các con số được công bố trên các benchmark phổ biến.

\subsection{Search-Based Software Testing (SBST) làm baseline}
\label{subsec:related-sbst}

Song song với hướng LLM, nhiều nghiên cứu dùng công cụ kiểm thử dựa trên tìm kiếm (SBST) như EvoSuite hoặc Pynguin làm đường cơ sở (baseline) để đối chiếu. Tang et al.~\cite{tang2024chatgptsbst} so sánh trực tiếp ChatGPT với EvoSuite trên 207 lớp Java, cho thấy EvoSuite vượt trội rõ rệt về statement coverage (74,2\% so với 55,4\%) và khả năng phát hiện lỗi ngoại lệ. Ryan et al.~\cite{ryan2024symprompt} (SymPrompt) dùng chính Pynguin làm baseline trên 897 phương thức Python -- các phương thức được chọn vì Pynguin \emph{không} đạt full coverage -- và ghi nhận Pynguin gốc đạt 72\% line coverage, cao hơn SymPrompt cơ bản (48\%) nhưng thấp hơn phiên bản đã lọc (77\%). Guerino và Vincenzi~\cite{guerino2025chatgptpython} cũng so sánh trực tiếp với Pynguin trên 40 chương trình Python, ghi nhận kết quả ngược lại: ChatGPT-3.5 (gộp nhiều temperature) vượt Pynguin cả về decision coverage (95,9\% so với 90,8\%) lẫn mutation score (91,5\% so với 69,8\%). Aminata et al.~\cite{aminata20252szsp} so sánh với EvoSuite trên 70 lớp Java, cho thấy EvoSuite thắng về coverage thô (56,1\% so với 43,4\%) nhưng thua về mutation score (30,8\% so với 41,3\% của LLM).

Các kết quả trên không thống nhất -- SBST có thể thắng hoặc thua LLM tuỳ dataset và metric -- nhưng đều gợi ý một điểm quan trọng: coverage cao không đảm bảo mutation score cao tương ứng, và ngược lại. Đáng chú ý, không nghiên cứu nào trong nhóm này xây dựng tập hàm được phân tầng có kiểm soát theo độ phức tạp cyclomatic (CC) với hạn ngạch cố định; SymPrompt và Guerino \& Vincenzi đều lấy hàm từ mã nguồn mở với phân bố CC không được kiểm soát, khiến độ phức tạp và hiệu năng mô hình có thể bị nhiễu lẫn nhau (confounded).

\subsection{Định vị nghiên cứu}
\label{subsec:related-positioning}

Khác với các nghiên cứu trên, đây là nghiên cứu đầu tiên đánh giá cụ thể \textbf{GPT-4o-mini} (thay vì GPT-4/GPT-4o đầy đủ) đối đầu trực tiếp với baseline \textbf{Pynguin} (thuật toán DYNAMOSA) trên một tập 50 hàm Python được phân tầng có kiểm soát theo ba mức độ phức tạp cyclomatic (CC 5--8, CC 9--12, CC 13--15) với hạn ngạch cố định 40\%/40\%/20\%, đồng thời đo lường song song cả độ bao phủ (branch coverage) lẫn khả năng phát hiện lỗi (mutation score) trên cùng một tập hàm cho cả hai phương pháp.
